\chapter{页故障}
\label{CH:PGFAULTS}
%% 

当使用了一个在页表中没有映射、或其映射的 \lstinline{PTE_V}
标志被清除、或其权限位（\lstinline{PTE_R}、
\lstinline{PTE_W}、
\lstinline{PTE_X}、
\lstinline{PTE_U}）
禁止了所尝试的操作的虚拟地址时，RISC-V CPU 会引发一个页故障异常。
RISC-V 区分三种
页故障：取数页故障（由取数指令引起）、
存数页故障（由存数指令引起）
和执行页故障（由要执行指令的取指引起）。
\lstinline{scause} 寄存器指示了页故障的类型，
而 \indexcode{stval} 寄存器则包含了无法转换的地址。

页表与页故障的组合是一个强大的工具。
页表为内核提供了一层虚拟地址到物理地址之间的间接性，使得内核能够控制地址空间的结构和内容。页故障则允许内核截获取数和存数操作，并通过修改页表来实时指定这些访问所引用的数据。内核可以利用这些能力来提高效率：例如，写时复制的 fork 使得内核能够在父进程与子进程之间透明地共享内存，从而避免复制双方都未写入的页面的开销。
应用程序员也能从中受益。其中一个可能性是内存映射文件，即内核利用分页机制使文件的内容出现在应用程序的地址空间中，并透明地根据页故障来读取文件块。另一个可能性是惰性内存分配，它允许程序申请一个巨大的虚拟地址空间，但只为程序实际读写的页面分配物理内存。xv6 仅将页故障用于一个目的：惰性分配。

在继续之前，请阅读 {\tt kernel/sysproc.c} 中的函数 {\tt sys\_sbrk()}，以及 {\tt kernel/vm.c} 中的 {\tt vmfault}。
在 {\tt kernel/trap.c} 和 {\tt kernel/vm.c} 中搜索对 {\tt vmfault} 的调用。

\section{惰性分配}
\label{sec:lazy}

xv6 的 \indextext{惰性分配} 包含两个部分。
首先，当应用程序通过调用带有 \lstinline{SBRK_LAZY} 标志的 \lstinline{sbrk} 请求内存时，内核会记录大小的增长，但不会为新范围的虚拟地址分配物理内存，也不会创建页表项（PTE）。其次，当这些新地址上发生页故障时，内核会分配一页物理内存并将其映射到页表中。内核实现的惰性分配对应用程序是透明的：应用程序无需修改即可从中受益。

惰性分配对应用程序来说非常方便，因为它们不需要精确预测自己将需要多少内存。例如，一个应用程序可能需要处理输入，但事先不知道输入会有多大。通过惰性分配，应用程序可以按最坏情况请求内存，但不必为这种最坏情况付出代价：对于应用程序从未使用的页面，内核完全不需要做任何工作。

此外，如果应用程序要求将地址空间增加很大的量，那么不使用惰性分配的 \lstinline{sbrk} 代价很高：如果应用程序请求 1 GB 的内存，内核必须分配并清零 262,144 个 4096 字节的物理页。惰性分配则可以将这种成本分摊到一段时间内。另一方面，惰性分配会带来额外的页故障开销，这涉及到用户/内核态的切换。操作系统可以通过每次页故障分配一批连续的页面（而不是一个页面），以及针对这类页故障专门优化内核的进入/退出代码（尽管 xv6 两者都没有做）来降低这种开销。

另一方面，当为惰性分配的页面发生页故障时，内核可能会发现没有空闲内存可供分配。在这种情况下，内核没有简单的方法向应用程序返回内存不足的错误，而是会杀死该应用程序。对于那些希望在分配失败时收到错误信息的应用程序，xv6 允许它们通过调用带有 \lstinline{SBRK_EAGER} 标志的 \lstinline{sbrk} 来急切地分配内存。

\section{代码}

系统调用 {\tt sbrk(n)} 会将进程的内存大小增加（如果 {\tt n} 为负数则减少）{\tt n} 字节，并返回新分配区域的起始地址（即旧的大小）。
内核的实现是 \lstinline{sys_sbrk}，见 \lineref{kernel/sysproc.c:/^sys_sbrk/}。

如果应用程序指定了 \lstinline{SBRK_EAGER}，该系统调用由函数 \lstinline{growproc} 实现，见 \lineref{kernel/proc.c:/^growproc/}。
\lstinline{growproc} 会调用 \lstinline{uvmalloc}。
\lstinline{uvmalloc}（见 \lineref{kernel/vm.c:/^uvmalloc/}）使用 {\tt kalloc} 分配物理内存，将分配的内存清零，然后使用 {\tt mappages} 在用户页表中添加页表项（PTE）。

如果应用程序进行惰性内存分配，\lstinline{sys_sbrk} 只是将进程的大小（\lstinline{myproc()->sz}）增加 {\tt n} 并返回旧的大小；它既不分配物理内存，也不向进程的页表中添加 PTE。

当一个进程加载或存储到某个缺少有效页表映射的虚拟地址时，CPU 会引发 \indextext{页故障异常}。
\lstinline{usertrap} 会检查这种情况（见 \lineref{kernel/trap.c:/page fault/}），并调用 \lstinline{vmfault}（见 \lineref{kernel/vm.c:/^vmfault/}）来处理页故障。\lstinline{vmfault} 会检查发生故障的地址是否在之前通过 {\tt sbrk} 授予的区域内，使用 {\tt kalloc} 分配一页物理内存，将分配的页面清零，并使用 {\tt mappages} 在用户页表中添加一个 PTE。Xv6 会在新页面的 PTE 中设置 \lstinline{PTE_W}、\lstinline{PTE_R}、\lstinline{PTE_U} 和 \lstinline{PTE_V} 标志。然后，\lstinline{usertrap} 会在引发故障的那条指令处恢复进程的运行。由于现在 PTE 已经有效，重新执行的取数或存数指令将不会再触发故障。

如果应用程序使用 {\tt sbrk} 释放内存，\lstinline{sys_sbrk} 会调用 \lstinline{shrinkproc}，后者会调用 \lstinline{uvmdealloc}。真正完成工作的是 {\tt uvmunmap}（见 \lineref{kernel/vm.c:/^uvmunmap/}），它使用 {\tt walk} 来查找 PTE。由于某些页面可能从未被进程使用过，因而也从未被 {\tt vmfault} 分配过，所以 {\tt uvmunmap} 会跳过那些没有 \lstinline{PTE_V} 标志的 PTE。如果某个 PTE 映射是有效的，{\tt uvmunmap} 会调用 {\tt kfree} 来释放它所指向的物理内存。

请注意，Xv6 使用进程的页表不仅是为了告诉硬件如何映射用户虚拟地址，也是作为记录该进程分配了哪些物理内存页的唯一依据。这就是为什么释放用户内存（在 {\tt uvmunmap} 中）需要检查用户页表。

\section{现实世界：写时复制（COW）fork}

许多内核（尽管不是 xv6）利用页故障来实现 \indextext{写时复制（COW）fork}。{\tt fork} 系统调用保证子进程在 fork 时刻看到的内存初始内容与父进程相同。一种实现方式是父进程的整个内存复制到子进程新分配的物理内存中；这正是 xv6 的做法。复制可能会很慢，如果子进程能够共享父进程的物理内存，效率会更高。然而，简单的共享实现是行不通的，因为父进程和子进程对共享的栈和堆的写入会相互干扰。

写时复制的 fork 通过适当使用页表权限和页故障，使得父进程和子进程能够安全地共享物理内存。基本计划是父进程和子进程初始时共享所有物理页，但每个进程都将这些页映射为只读（即清除 \lstinline{PTE_W} 标志）。这样父进程和子进程都可以从共享的物理内存中读取。如果其中任何一个进程写入某个共享页，RISC-V CPU 就会引发一个页故障异常。支持 COW 的内核会通过分配一页新的物理内存，并将共享页的内容复制到这个新页中来响应。然后，内核会将引发故障的进程的页表中的相应 PTE 修改为指向这个副本，并同时允许读写，最后在引发故障的那条指令处恢复该进程的执行。由于现在 PTE 允许写入，重新执行的存数指令将不再触发故障，并且修改的是页面的私有副本，而非共享页面。

写时复制需要额外的记账信息，以帮助决定何时可以释放物理页，因为根据 fork、页故障、exec 和 exit 的历史记录，每个物理页可能被不同数量的页表所引用。这种记账机制还允许一个重要的优化：如果一个进程发生了存储页故障，并且该物理页只被该进程的页表所引用，那么就不需要进行复制。

写时复制使得 \lstinline{fork} 更快，因为 \lstinline{fork} 不需要复制内存。某些内存可能会在之后被写入时才需要复制，但通常大部分内存永远都不需要复制。一个常见的例子是 \lstinline{fork} 后跟着 \lstinline{exec}：在 \lstinline{fork} 之后可能只有少数几个页面会被写入，然后子进程的 \lstinline{exec} 会释放从父进程继承来的大部分内存。写时复制的 \lstinline{fork} 完全省去了复制这些内存的必要。此外，COW fork 是透明的：应用程序无需任何修改即可从中受益。

\section{现实世界：请求分页}

另一个广泛使用页故障的特性是 \indextext{请求分页}。在 \lstinline{exec} 系统调用中，xv6 会在启动应用程序之前将其所有的文本和数据都加载到内存中。由于应用程序可能很大，而且从磁盘读取需要时间，这种启动开销对用户来说可能是明显的。为了减少启动时间，现代内核不会一开始就将可执行文件加载到内存中，而是只创建用户页表，并将所有 PTE 标记为无效。内核随后启动程序运行；每当程序第一次使用某个页面时，就会发生一次页故障，作为响应，内核会从磁盘读取该页面的内容，并将其映射到用户的地址空间中。与 COW fork 和惰性分配一样，内核可以以对应用程序透明的方式实现这一特性。

计算机上运行的程序可能需要的内存可能超过计算机拥有的 RAM。为了优雅地应对这种情况，操作系统可能会实现 \indextext{分页到磁盘}。其思路是只将用户页面的一部分存放在 RAM 中，而将其余部分存放在磁盘的 \indextext{分页区域} 中。内核将对应于存放在分页区域（因此不在 RAM 中）的内存的 PTE 标记为无效。如果某个应用程序试图使用一个已被{\it 换出}到磁盘的页面，应用程序将触发一个页故障，然后该页面必须被{\it 换入}：内核的陷阱处理程序会分配一页物理 RAM，从磁盘将该页面读到 RAM 中，并修改相应的 PTE 使其指向该 RAM。

当需要换入一个页面，但没有空闲的物理 RAM 时会发生什么？在这种情况下，内核必须首先通过将一个物理页换出（或{\it 驱逐}）到磁盘上的分页区域，并标记指向该物理页的 PTE 为无效，来释放一个物理页。驱逐操作代价高昂，因此分页在不频繁发生时性能最佳：即应用程序只使用其内存页的一个子集，并且这些子集的并集能够放入 RAM 中。这一性质通常被称为具有良好的引用局部性。与许多虚拟内存技术一样，内核通常以对应用程序透明的方式实现分页到磁盘。

无论硬件提供多少 RAM，计算机通常都只有很少甚至没有{\it 空闲}物理内存。例如，云提供商为了经济高效地利用硬件，会在单台机器上复用大量客户。又例如，用户在有限的物理内存中的智能手机上运行许多应用程序。在这种环境下，分配一个页面可能首先需要驱逐一个现有页面。因此，当空闲物理内存稀缺时，分配操作的成本很高。

在空闲内存稀缺且程序只活跃地使用其已分配内存的一小部分时，惰性分配和请求分页尤其有利。这些技术还可以避免当页面被分配或加载后，要么从未使用，要么在使用之前就被换出时浪费的工作。

\section{现实世界：内存映射文件}

其他结合了分页和页故障异常的特性包括自动增长的栈和 \indextext{内存映射文件}。所谓内存映射文件，是指程序使用 \texttt{mmap} 系统调用将其映射到自己的地址空间中，从而可以通过取数和存数指令来读写这些文件。

% “虚拟”内存，换出，换入
% 惰性分配
% 自动栈扩展
% 保护页
% mmap 文件
% cow fork
% 共享文本，共享库
% 请求分页文本
% 虚拟机迁移
% 分布式共享内存
% 快速 IPC
% 零拷贝写入
% DPDK
% （统一块 / 页缓存）

%% 
\section{练习}
%% 

\begin{enumerate}

\item 编写一个用户程序，通过调用 \lstinline{sbrk(1)} 将其地址空间增加一个字节。
运行该程序，并在调用 \lstinline{sbrk} 之前和之后检查程序的页表。
内核分配了多少空间？新内存对应的 PTE 中包含什么内容？

\item 实现 COW fork。

\item 实现 {\tt mmap}。

\end{enumerate}
